\chapter{Background: $\infty$-cosmos}
\label{chap:infc}
\section{$\infty$-cosmos}
\label{sec:infct}
The definition of an infinity-cosmos was introduced by Riehl and Verity over the course of a series of papers. That definition has evolved over time. We start by reviewing the current form of the definition from \cite{RV-elements} first, as it is the simplest definition. But then we also mention a more general older definition, which will be the one that is relevant to the project here. For the sake of keeping naming of these definitions clear, I will call the current form as the $\infty$-cosmos, the older more general form as Non-cofibrant $\infty$-cosmos, and the one with my modification as weak $\infty$-cosmos.

\begin{defn}[ $\infty$-cosmos]
    A \emph{strict $\infty$-cosmos $\mathcal{K}$} is a category that is enriched over quasi-categories, where given two objects $A, B$, we denote the quasi-category of morphisms as $\cat{Fun(A,B)}$ and refer to it as a \emph{functor space}, that is also equipped with a specified collection of maps that we call \emph{isofibrations} and denote by "$\twoheadrightarrow$" satisfying the following two axioms:
    \begin{enumerate}
        \item (completeness) The quasi-categorically enriched category $\mathcal{K}$ possesses a terminal object, small products, pullbacks of isofibrations, limits of countable towers of isofibrations, and cotensors with simplicial sets, each of these limit notions satisfying a universal property that is enriched over simplicial sets.
        \item (isofibrations) The isofibrations contain all isomorphisms and any map whose codomain is the terminal object; are closed under composition, product, pullback, forming inverse limits of towers, and Leibniz cotensors with monomorphism of simplicial sets; and have the property that if $f:A\twoheadrightarrow B$ is an isofibration and $X$ is any object then $\cat{Fun(X,A)}\twoheadrightarrow \cat{Fun(X,B)}$ is an isofibration of quasi-categories.
    \end{enumerate}
\end{defn}

The axioms are inspired by the properties of the fibrant objects in a model structure that
\begin{enumerate}
    \item is enriched over the Joyal's model structure on quasicategories and
    \item in which all fibrant objects are also cofibrant.
\end{enumerate}
 The more general form of this definition takes away this second assumption. In terms of the difference in the mathematical objects, you can prove the same with both definitions, but the proofs may be more stream lined with the additional assumption (thus the naming convention). 

The more general definition come from the paper \cite{rvp} 

\begin{defn}[Non-Cofibrant $\infty$-Cosmos]
     A \emph{strict $\infty$-cosmos $\mathcal{K}$} is a simplicially enriched category together with two specified classes of its $0$-arrows $\mathcal{W}$ called its class of \textit{weak equivalences}, $\mathcal{F}$ called its class of \textit{isofibrations} and all maps in $\mathcal{W} \cap \mathcal{F}$ the \textit{trivial fibrations.} These classes satisfythe following axioms:
     \begin{enumerate}
        \item (completeness) the category $\mathcal{K}$ possesses a terminal object, cotensors with all finitely presented simplicial sets, and pullbacks of isofibrations along any morphism;
        \item (fibrant objects) all of the morphisms with codomain the terminal object are isofibrations;
        \item (pullback stability) the classes of isofibrations and trivial fibrations are both stable under pullback along all morphisms;
        \item (Leibniz stability) if $p:E\twoheadrightarrow B$ is an isofibration in $\mathcal{K}$ and $i:U\hookrightarrow V$ is an inclusion of finitely presented simplicial sets then the Leibniz cotensor is an isofibration and it is a trivial fibration when $p$ is a trivial fibration in $\mathcal{K}$ or $i$ is trivial cofibration in the Joyal model structure on $\cat{sSet}$; and
        \item (cofibrant replacement) if $A$ is an arbitrary object of $\mathcal{K}$ then there exist a trivial fibration $r_A:\bar{A}\stackrel{\sim}{\twoheadrightarrow} A$ for which the object $\bar{A}$ is \emph{cofibrant} in the sense that enjoys the left lifting property with respect to all trivial fibrations in $\mathcal{K}$.
     \end{enumerate}
     %, equipped with two specified classes of its $0$-arrows $\mathcal{W}$ called its class of \textit{weak equivalences} and $\mathcal{F}$ called its class of \textit{isofibrations}. As usual, we shall assume that $\mathcal{W}$ contains all isomorphisms and satisfies the \textit{$2$-of-$6$ porperty} and that $\mathcal{F}$ also contains all isomorphisms and is closed under composition. We use the notation $\cat{Map(A,B)}$ to denote the hom-space between the objects $A$ and $B$ of $\mathcal{K}$.
\end{defn}
\begin{rmk}\label{remarkabout}
    Every $\bic \ \mathcal{K}$ is a $\pic$ where the weak equivalences are defined by the following:$f : A \rightarrow B$ is a weak equivalence when the induced map $f_*: \cat{Fun(X,A)}\stackrel{\sim}{\rightarrow} \cat{Fun(X,B)}$ on functor spaces is an equivalence of quasi-categories for all $X\in \mathcal{K}$. See more in \cite{RV-elements} $\S~1.2$
    
\end{rmk}

We can see in this more general form the model structure working in the background. In fact a lot of examples of $\infty$-cosmoi arising in nature come from "quasi-categorically enriched categories of fibrant objects". What this mean is a combination of the following two definitions:

\begin{defn}[category of fibrant objects]
    A \textbf{category of fibrant objects} consists of a category $\mathcal{M}$ together with two subcategories of morphisms $\cat{W}$ and $\cat{F}$ satisfying the following conditions:
    \begin{enumerate}
        \item $\mathcal{M}$ has products and in particular a terminal object $\cat{1}$. Moreover, the classes $\cat{F}$ and $\cat{F}\cap\cat{W}$ are each closed under products.
        \item Every morphism in $\cat{W}$ satisfies the $2$-of-$3$ property: if two out of three morphisms $f$, $g$, and $g\circ f$ are in $\cat{W}$, then so is the third.
        \item Pullbacks of maps in $\cat{F}$ along any morphism exist and are again in $\cat{F}$, and the class $\cat{F}\cap\cat{W}$ is also stable under pullback.
        \item Limits of countable towers of maps in $\cat{F}$ exist and are again in $\cat{F}$. Likewise for maps in $\cat{F}\cap\cat{W}$.
        \item For every object $B$, there exists a \textbf{path object} $PB$ together with a factorization of the diagonal into a map in $\cat{W}$ followed by a map in $\cat{F}$:
        \[
        \begin{tikzcd}
            {} & PB \arrow[dr, two heads] & \\
            B \arrow[ur, "\sim" {sloped}] \arrow[rr, "\Delta"'] & {} & B \times B
        \end{tikzcd}
        \]
        \item All objects are fibrant: for every $B\in \mathcal{M}$, the map $B \to \cat{1}$ lies in $\cat{F}$.
    \end{enumerate} 
\end{defn}
\begin{defn}[Quasicategorically enriched model structure]
    Let $\mathcal{V}:=\cat{sSet}$, the cartesian closed model category, using the Joyal model structure. A \textbf{$\mathcal{V}$-model category} is a model category $\mathcal{M}$ that is tensored, cotensored, and $\mathcal{V}$-enriched and so that
    \begin{enumerate}
        \item $(\otimes,\{,\},\hom)$ is a Quillen two-variable adjunction and
        \item the map $Q1\otimes m \rightarrow \cat{1}\otimes m \equiv m$ defined by the cofibrant replacement of the terminal object is a weak equivalence whenever $m$ is cofibrant.
    \end{enumerate}
\end{defn}

Given the fact that the axioms are inspired by the conditions previously written, it is not surprising the fact many naturally occurring $\infty$-cosmoi are realized as quasi-categorically enriched categories of fibrant objects. Then the following theorem found in \cite{RV-elements} provides a concrete realization of this idea.

\begin{thm}
   Let $\mathcal{M}$ be any model category that is enriched over the Joyal model structure and in which every fibrant object is cofibrant. Then the full subcategory of fibrant objects $\mathcal{M}_f$ inherits the structure of an $\bic$ in which the isofibrations are the fibrations between fibrant objects, the equivalences are the weak equivalences between fibrant objects, and the trivial fibrations are the trivial fibrations between fibrant objects.
\end{thm}

I will write a sketch of the proof for the slightly more general version of the proposition where we drop the second assumption that every fibrant object is cofibrant.

\begin{prop}
    Let $\mathcal{M}$ be any model category that is enriched over the Joyal model structure. Then the full subcategory of fibrant objects $\mathcal{M}_f$ inherits the structure of an $\pic$ in which the isofibrations are the fibrations between fibrant objects, the equivalences are the weak equivalences between fibrant objects, and the trivial fibrations are the trivial fibrations between fibrant objects.
\end{prop}

\begin{proof}
    Let $\mathcal{M}$ be a model category enriched over the Joyal model structure. That means that is enriched, tensored, and cotensored over $\sset_{Joyal}$. Let $\mathcal{M}_f\subseteq \mathcal{M}$ be the full simplicially enriched subcategory spanned by the fibrant objects.
    Define:
    \[\mathcal{W}_{\mathcal{M}_f}=\left\{f:X\rightarrow Y \| X,Y \in \mathcal{M}_f, f \text{ is a weak equivalence in }\mathcal{M}\right\}\]
    \[\mathcal{F}_{\mathcal{M}_f}=\left\{p:E\rightarrow B \| E,B \in \mathcal{M}_f, p \text{ is a fibration in }\mathcal{M}\right\}\]
    
    We verify the axioms for $\pic$. We get most of the completeness (except for cotensors with simplicial sets), fibrant objects, pullback stability, and cofibrant replacements given by the fact that $\mathcal{M}_f$ is the full category generated by the fibrant objects of $\mathcal{M}$. 
   
   The ones missing are those involving the interaction between $\mathcal{M}_f$ and simplicial sets. For these, we use the enrichment over the Joyal model structure. This enrichment gives us the followings functors

   \[\hom:\mathcal{M}^\op\times \mathcal{M}\rightarrow\sset \]
   \[\otimes:\mathcal{M}\times \sset \rightarrow \mathcal{M}\]
   \[\{,\}:\sset^\op\times \mathcal{M} \rightarrow \mathcal{M}\]
   
   that satisfy the following criteria:
   \begin{itemize}
    \item For $X,Y \in \mathcal{M}$ and $U \in \sset$, we have the following isomorphisms in $\sset$ \[\hom(X,Y)^U\cong\hom(X\otimes U,Y)\cong \hom(X,\{U,Y\})\]
    \item (Pullback-Power axiom) If $i:U\hookrightarrow V$ is a cofibration in $\sset_{Joyal}$ and $p:X\to Y$ is a fibration between fibrant objects in $\mathcal{M}$, then the induced map
   \[
        X^V \to Y^V \times_{Y^U} X^U
   \]
   is a fibration between fibrant objects, and it is a trivial fibration if either $i$ or $p$ is a trivial cofibration.
   \end{itemize}
    We can see that the leibniz stability axiom is satisfied by the second condition. 
\end{proof}

\begin{rmk}
We now note the key structural differences between $\bic$ and $pic$. In the $\bic$ definition, the assumption of cofibrancy ensures that the functor spaces $\cat{Fun(A,B)}$ are quasi-categories, giving an enrichment over the subcategories of $\qcat$. When the cofibrant assumption is dropped for $\pic$, we no longer can guarantee this on the nose. To compensate for this loss property we must require some more structure on the definition of $\bic$. This comes from the fact that we require to define the family of weak equivalences as additional data instead of being given by the functor space (see remark \ref{remarkabout}) and the axiom of cofibrant replacement.
\end{rmk}

%But the way that I would actually write it would be out of order: write the definitions first and then figure out what needs to be explained and then write that before the definitions in the paper.

%Next thing would be to explain how you use a model category enriched over Joyal model structure to get examples. Appendix E
%Section 3.2 is "Examples of oo-cosmoi" which contains an exposition of section E.1 but adapted to the more general setting.